% Lösungen Einheit 2
\einheitenkopf{Einheit 2 von 5 · Extrempunkte}
\erg{14}{a) gesucht \quad b) gegeben \quad c) gegeben \quad d) gesucht \quad e) gegeben}
\erg{15}{a) A \quad b) B \quad c) D \quad d) C \quad e) B (am Faktor $2 - x$ ablesbar; A ginge auch, ist aber aufwendiger)}
\erg{16}{a) $f'(x) = x^2 - 4$: $x = 2$, $x = -2$ \quad b) $f'(x) = x^2 - 2x - 3$: $x = 3$, $x = -1$ \quad c) $f'(x) = 6x^2 + 6x - 12$, $x^2 + x - 2 = 0$: $x = 1$, $x = -2$ \quad d) $f'(x) = -x^2 + 4x + 5$: $x = 5$, $x = -1$}
\erg{17}{a) $f''(x) = 6x + 6$; $H(-3 \mid 27)$, $T(1 \mid -5)$ \quad b) $f'(x) = 4x(x - 1)(x - 2)$; $T(0 \mid 0)$, $H(1 \mid 1)$, $T(2 \mid 0)$ \quad c) $f'(x) = 12x^2(x - 1)$; $T(1 \mid -1)$; bei $x = 0$: $f''(0) = 0$, $f'''(0) = -24 \ne 0$, $f'$ ohne Vorzeichenwechsel: Sattelpunkt $S(0 \mid 0)$ \quad d) $f'(x) = x^3 - 7x + 6 = (x - 1)(x^2 + x - 6)$; Extremstellen $1$, $2$, $-3$; $f''(x) = 3x^2 - 7$: Hochpunkt bei $x = 1$, Tiefpunkte bei $x = 2$ und $x = -3$ \quad e) $f'(x) = (x^2 + 2x - 3)\,\mathrm{e}^{x}$; $H(-3 \mid 6\mathrm{e}^{-3}) \approx (-3 \mid 0{,}30)$, $T(1 \mid -2\mathrm{e}) \approx (1 \mid -5{,}44)$ \quad f) $a'(t) = 5\,\mathrm{e}^{-0{,}2t}(1 - 0{,}2t) = 0 \Leftrightarrow t = 5$; $a(5) = 25\,\mathrm{e}^{-1} \approx 9{,}2$: nach 5 Minuten etwa 9 Anrufe pro Minute \quad g) $g'(x) = -3x(x - 4)$; $T(0 \mid 0)$, $H(4 \mid 32)$; $g(-3) = 81$, $g(5) = 25$: größter Wert $81$ (bei $x = -3$, am Rand), kleinster Wert $0$ (bei $x = 0$) \quad h) $f'(x) = 4x(x^2 - 1)$; $T(\pm 1 \mid 2)$, $H(0 \mid 3)$; $f(x) \to \infty$: Wertebereich $[2;\,\infty[$ \quad i) $f'(x) = (2x - 0{,}5x^2)\,\mathrm{e}^{-0{,}5x} = 0{,}5x\,(4 - x)\,\mathrm{e}^{-0{,}5x}$; Vorzeichenwechsel bei $0$ von $-$ nach $+$: $T(0 \mid 0)$; bei $4$ von $+$ nach $-$: $H(4 \mid 16\mathrm{e}^{-2}) \approx (4 \mid 2{,}17)$}
\erg{18}{a) $f'(5) = 2\cdot 5 - 10 = 0$ \quad b) $f'(3) = 27 - 27 = 0$ \quad c) $f'(2) = 24 + 12 - 36 = 0$ \quad d) $f'(-3) = -108 + 108 = 0$ \quad e) $f'(4) = -48 + 24 + 24 = 0$, $f''(4) = -18 < 0$: Hochpunkt \quad f) $f'(-1) = -4 + 4 = 0$, $f''(-1) = 12 > 0$, $f(-1) = -7$ \quad g) $f'(x) = 4(x + 2)^3$: für $x < -2$ negativ, für $x > -2$ positiv, Wechsel von $-$ nach $+$: Tiefpunkt $(-2 \mid -1)$ \quad h) $f'(3) = 0$ und $f''(3) < 0$ ist die hinreichende Bedingung für einen Hochpunkt \quad i) $f'(x) = -3(x - 1)^2$: $f'(1) = 0$, aber $f'(x) \le 0$ auf beiden Seiten, kein Vorzeichenwechsel: Sattelpunkt $S(1 \mid 3)$ \quad j) $f'(0) = -3 < 0$, $f'(1) = 2 > 0$: $f'$ wechselt dazwischen von $-$ nach $+$, also liegt dort ein Tiefpunkt (bei $x \approx 0{,}72$) \quad k) $f'(x) = 12x^2(x - 2)$, $f'(2) = 0$, $f''(2) = 48 > 0$: $T(2 \mid -14)$; $f'(0) = 0$, aber $12x^2 \ge 0$ und $x - 2 < 0$ nahe $0$, also $f' \le 0$ auf beiden Seiten: kein Vorzeichenwechsel, kein Extrempunkt}
\erg{19}{a) Nullstellen $1$ und $5$, Extremstelle $3$ \quad b) Scheitel $S(2 \mid 3)$ ist der tiefste Punkt; kleinster Abstand $3 - 1 = 2$ \quad c) $\mathrm{e}^{-x} \ne 0$, also nur $x^2 - 4 = 0$: $x = \pm 2$; $f'(x) = -(x^2 - 2x - 4)\,\mathrm{e}^{-x} = 0 \Leftrightarrow x = 1 \pm \sqrt{5}$; Tiefpunkt bei $x = 1 - \sqrt{5} \approx -1{,}24$ (Vorzeichenwechsel von $-$ nach $+$) \quad d) $f'(x) = 0{,}75x^2 - 3 = 0 \Leftrightarrow x = \pm 2$; $T(2 \mid -4)$, $H(-2 \mid 4)$; Abstand $\sqrt{4^2 + 8^2} = \sqrt{80} \approx 8{,}94$ \quad e) $f'(x) = -1{,}5x^2 + 1{,}5 = 0 \Leftrightarrow x = \pm 1$; $f''(x) = -3x$: $H(1 \mid 1)$, $T(-1 \mid -1)$; beide Punkte erfüllen $y = x$}
\erg{20}{a) Jonas vertauscht die Art: $f''(-1) < 0$ bedeutet Hochpunkt, $f''(3) > 0$ Tiefpunkt; richtig $H(-1 \mid 7)$, $T(3 \mid -25)$ \quad b) $g'(x) = 3x(x + 2)$, $g''(x) = 6x + 6$: $H(-2 \mid 3)$, $T(0 \mid -1)$}
\erg{21}{a) Auch an einer Sattelstelle ist $f'(x_0) = 0$, z. B. $f(x) = x^5$ bei $0$: $f'$ wechselt dort das Vorzeichen nicht. \quad b) Der größte Wert kann am Rand liegen: steigt $f$ bis $b$, ist $f(b)$ am größten, obwohl $f'(b) \ne 0$. \quad c) stimmt: $f'$ ist quadratisch mit zwei einfachen Nullstellen und wechselt an beiden das Vorzeichen \quad d) falsch: $f'(x) > 0$ stimmt, aber $f(x) > y_0$, weil $y_0$ in der Nähe der kleinste Wert ist}
\erg{22}{a) Schnittpunkte $(0 \mid 0)$ und $(3 \mid 0)$; $f'(-1) = 9 > 0$, $f'(1) = -3 < 0$: Vorzeichenwechsel von $+$ nach $-$ bei $0$, Hochpunkt $H(0 \mid 0)$ \quad b) Der Hochpunkt liegt auf der $x$-Achse; in seiner Nähe ist $f(x) \le 0$, der Graph bleibt auf einer Seite: Berührpunkt \quad c) $f'(x) = \mathrm{e}^{x} - 1 = 0 \Leftrightarrow x = 0$, $f''(0) = 1 > 0$: $T(0 \mid 1)$; $f'$ hat nur diese Nullstelle, also keine weiteren Extrempunkte \quad d) $f'(x) = -2x\,\mathrm{e}^{-x^2}$: $f$ steigt für $x \le 0$ und fällt für $x \ge 0$, also größter Wert $f(0) = 1$; $f(x) > 0$ und $f(x) \to 0$ für $x \to \pm\infty$, Graph symmetrisch zur $y$-Achse; Wertemenge $]0;\,1]$}
\erg{23}{a) $f'(x) = -x + 2 = 0 \Leftrightarrow x = 2$, Gipfel $(2 \mid 3)$; $g'(x) = -0{,}75x^2 + 7{,}5x - 18 = -0{,}75(x - 4)(x - 6)$, $g''(6) = -1{,}5 < 0$: Gipfel $(6 \mid 2)$ \quad b) Abstand $\sqrt{4^2 + 1^2} = \sqrt{17} \approx 4{,}12$, also etwa 412 m $\le$ 450 m: das Seil ist möglich}
